%% ============================================================================
%% AURORA — ACM SIGPLAN Proceedings IMRAD Paper (Overleaf Ready)
%% Template: Association for Computing Machinery (ACM) - SIGPLAN Proceedings
%% Overleaf Project Link: https://www.overleaf.com/project/6aeaa2e64e1d2f7cf505e2cb
%% Instructions:
%%   1. Paste this entire file into `main.tex` (or `sample-sigplan.tex`) in Overleaf.
%%   2. Paste `references.bib` into your Overleaf bibliography file.
%%   3. Click "Recompile" in Overleaf to generate the official 2-column ACM PDF.
%% ============================================================================
\documentclass[sigplan,screen,nonacm]{acmart}

\usepackage{booktabs}
\usepackage{multirow}
\usepackage{amsmath}
\usepackage{listings}
\usepackage{xcolor}

%% Remove ACM copyright / conference footer clutter for institutional capstone submission
\settopmatter{printacmref=false}
\renewcommand\footnotetextcopyrightpermission[1]{}
\pagestyle{plain}

\begin{document}

%% ============================================================================
%% TITLE & AUTHORS
%% ============================================================================
\title[AURORA: Paperless Academic Defense Workflow & QA Platform]{AURORA: An Academic Unified Review, Observation, Rating, and Assessment Platform for Paperless Capstone Defense Workflows}
\subtitle{Department of Computational Sciences, College of Engineering and Computational Sciences, Partido State University}

\author{Gene Vincent Ma\~{n}ago}
\affiliation{%
  \institution{Partido State University}
  \department{Dept. of Computational Sciences}
  \city{Goa, Camarines Sur}
  \country{Philippines}
}
\email{gmanago@parsu.edu.ph}

\author{Carol Ann Castillo}
\affiliation{%
  \institution{Partido State University}
  \department{Dept. of Computational Sciences}
  \city{Goa, Camarines Sur}
  \country{Philippines}
}
\email{ccastillo@parsu.edu.ph}

\author{Rona Mae Buenafe}
\affiliation{%
  \institution{Partido State University}
  \department{Dept. of Computational Sciences}
  \city{Goa, Camarines Sur}
  \country{Philippines}
}
\email{rbuenafe@parsu.edu.ph}

\author{Lara Mae Encinas}
\affiliation{%
  \institution{Partido State University}
  \department{Dept. of Computational Sciences}
  \city{Goa, Camarines Sur}
  \country{Philippines}
}
\email{lencinas@parsu.edu.ph}

\author{Froiland B. Pe\~{n}alosa}
\authornote{Research Adviser, Department of Computational Sciences, Partido State University.}
\affiliation{%
  \institution{Partido State University}
  \department{Dept. of Computational Sciences}
  \city{Goa, Camarines Sur}
  \country{Philippines}
}
\email{fpenalosa@parsu.edu.ph}

\renewcommand{\shortauthors}{Ma\~{n}ago, Castillo, Buenafe, Encinas, and Pe\~{n}alosa}

%% ============================================================================
%% ABSTRACT (IMRAD SUMMARY)
%% ============================================================================
\begin{abstract}
Traditional undergraduate capstone and thesis defense workflows in Philippine state universities rely heavily on physical manuscript printing, manual adviser endorsement routing, ad-hoc room scheduling, and hand-computed paper score sheets. These fragmented processes result in prolonged review turnaround times, scheduling collisions, high financial printing costs for students, and the absence of cryptographic auditability. This paper presents \textbf{AURORA} (\textit{Academic Unified Review, Observation, Rating, and Assessment}), a cloud-native, paperless academic defense management and quality assurance platform engineered for the College of Engineering and Computational Sciences (CECS) at Partido State University. Built using Next.js 16 (App Router), React 19, TypeScript, and Supabase (PostgreSQL 15 with Row-Level Security), AURORA unifies four core architectural innovations: (1) a resolution-independent, dual-pane \texttt{PDF.js} manuscript review workspace with real-time WebSocket coordinate annotations and FIPS 180-4 SHA-256 document version fingerprinting; (2) an automated defense scheduler with venue clash detection and an academic Conflict-of-Interest Firewall that prevents research advisers from sitting as voting panelists for their own advisees; (3) a dynamic weighted rubric engine (Form DCS-CF-04) integrated with Republic Act No. 8792 step-up password re-authentication and tamper-evident digital signature serialization (\texttt{AURORA-YYYY-NNNNNN}); and (4) an administrative self-healing maintenance and 3-2-1 disaster recovery console. Empirical evaluation with $N = 27$ institutional stakeholders (students, advisers, panelists, capstone coordinator, and college dean) using a validated 10-item System Usability Scale (SUS) and 27-item Unified Theory of Acceptance and Use of Technology (UTAUT) instrument ($\alpha = 0.946$) yielded a composite SUS score of \textbf{91.8 / 100 (Grade A --- Excellent)} and a grand UTAUT mean of \textbf{4.86 / 5.00 (Highly Acceptable)}, while reducing manuscript endorsement turnaround time by \textbf{85.9\%} (from 18.4 days to 2.6 days).
\end{abstract}

\keywords{Paperless Academic Workflow, Capstone Defense Management, Coordinate PDF Annotation, Republic Act No. 8792 Digital Signatures, Row-Level Security, System Usability Scale (SUS), UTAUT}

\maketitle

%% ============================================================================
%% 1. INTRODUCTION
%% ============================================================================
\section{Introduction}

Undergraduate capstone research is a cornerstone requirement for computing degrees under Commission on Higher Education (CHED) Memorandum Order No. 25, Series of 2015 \cite{ched2015cmo25}. At Partido State University (ParSU), students enrolled in the Bachelor of Science in Information Technology (BSIT) program under the College of Engineering and Computational Sciences (CECS) progress through a rigorous multi-stage oral defense lifecycle encompassing Concept Paper Defense, Title Defense, Proposal Defense, and Final Oral Defense.

Despite the digital nature of computing education, the operational administration of capstone defenses has historically remained bound to physical paper workflows. Student teams print multiple copies of 80-to-150-page manuscripts for every defense iteration; research advisers annotate drafts manually with ink; defense coordinators cross-reference faculty schedules on spreadsheets to book deliberation rooms; and panel evaluators compute weighted grades manually on printed Form DCS-CF-04 sheets. Key informant interviews conducted at ParSU CECS revealed four systemic bottlenecks in this traditional paradigm:
\begin{enumerate}
  \item \textbf{Prolonged Review Turnaround:} Physical routing of manuscripts and Form DCS-CF-03 endorsement slips requires an average of 18.4 calendar days per cycle.
  \item \textbf{Scheduling & Impartiality Risks:} Manual scheduling is prone to double-booked venues, overlapping panelist timeslots, and inadvertent conflict-of-interest assignments.
  \item \textbf{Grading & Archival Fragility:} Hand-computed rubric sheets are susceptible to arithmetic errors, misplaced physical forms, and unauthorized post-deliberation alterations.
  \item \textbf{Operational & Disaster Vulnerability:} Decentralized paper and spreadsheet records lack automated integrity verification, self-healing maintenance, and institutional disaster recovery backups.
\end{enumerate}

\subsection{Research Objectives}
To address these institutional challenges, this study designed, implemented, and evaluated the \textbf{AURORA} platform in accordance with four specific objectives:
\begin{itemize}
  \item \textbf{Objective 1 (Workflow Digitization \& Gating):} Automate the multi-stage capstone defense lifecycle---including college-customizable defense stages, SHA-256 manuscript version control, and Form DCS-CF-03 adviser endorsement gating.
  \item \textbf{Objective 2 (Real-Time Collaborative Review \& Conflict-Free Scheduling):} Engineer a split-screen interactive PDF annotation workspace synchronized via WebSockets alongside an automated defense scheduler with room-clash and adviser conflict-of-interest guards.
  \item \textbf{Objective 3 (Cryptographic Rubric Deliberation \& Maintainability):} Implement a dynamic weighted rubric engine (Form DCS-CF-04) secured by Republic Act No. 8792 step-up digital signatures, public certificate verification, and a 3-2-1 Disaster Recovery and Self-Healing Maintenance Console.
  \item \textbf{Objective 4 (Usability \& Technology Acceptance Evaluation):} Quantitatively evaluate the system's usability and acceptability among institutional stakeholders using the System Usability Scale (SUS) \cite{brooke1996sus} and the Unified Theory of Acceptance and Use of Technology (UTAUT) \cite{venkatesh2003user}.
\end{itemize}

%% ============================================================================
%% 2. RELATED WORK
%% ============================================================================
\section{Related Work}

\subsection{Academic Manuscript & Defense Management Systems}
Learning Management Systems (LMS) such as Moodle and Google Classroom facilitate generic file submissions but lack specialized academic governance primitives required for oral defenses: multi-stage prerequisite gating, panel committee composition rules, weighted multi-evaluator rubric consensus, and verifiable clearance serials. Prior university repository systems focus primarily on post-graduation archival rather than active, iterative pre-defense review and live deliberation scoring.

\subsection{Viewport-Independent PDF Annotation}
Conventional document review tools require downloading files locally, breaking version continuity. Browser-native vector rendering via \texttt{PDF.js} enables zero-install inspection; however, pixel-based annotation coordinates drift when viewed across varying monitor resolutions or zoom levels. Storing bounding boxes as normalized floating-point percentages $(x\%, y\%, w\%, h\%) \in [0, 100]$ ensures mathematical invariance across desktop, laptop, and tablet viewports.

\subsection{Electronic Signatures & Legal Non-Repudiation}
Under Philippine Republic Act No. 8792 (\textit{Electronic Commerce Act of 2000}) \cite{ra8792}, an electronic signature possesses legal admissibility only when a reliable procedure authenticates the signer's identity and prevents subsequent alteration of the signed payload. Combining step-up credential re-authentication with FIPS 180-4 SHA-256 cryptographic hashing \cite{nist2015fips180} and PostgreSQL Row-Level Security (RLS) immutability policies satisfies statutory non-repudiation requirements.

%% ============================================================================
%% 3. METHODOLOGY (MATERIALS AND METHODS)
%% ============================================================================
\section{Methodology}

\subsection{Software Development Life Cycle (Agile Scrum)}
The development of AURORA followed the \textbf{Agile Scrum Software Development Life Cycle (SDLC)} across a 14-month academic timeline spanning Academic Year 2025--2026 (1st Semester, 2nd Semester, and Midyear Term) through the 1st Semester of Academic Year 2026--2027 (August 15, 2025 to October 2, 2026). Ten iterative milestones (M-01 to M-10) delivered incremental, testable software modules validated continuously with the Research Adviser and Capstone Coordinator.

\subsection{Four-Tier Cloud System Architecture}
AURORA is architected as a decoupled four-tier cloud-native web application (Table~\ref{tab:architecture}):
\begin{enumerate}
  \item \textbf{Client Presentation Tier:} Built with React 19 and Tailwind CSS v4, featuring a 60/40 dual-pane review workspace (\texttt{PDF.js} vector canvas + rubric drawer) and an HTML5 Canvas electronic signature pad.
  \item \textbf{Edge \& Server Action Tier:} Hosted on the Vercel Global Edge Network using Next.js 16 (App Router). Privileged mutations execute exclusively within type-safe Server Actions protected by Edge Middleware RBAC route guards.
  \item \textbf{Backend Service Tier:} Powered by Supabase Cloud (AWS \texttt{ap-southeast-1} Singapore), providing GoTrue HTTP-only JWT authentication, private S3 storage buckets (\texttt{/manuscripts}, \texttt{/signatures}), and Realtime WebSocket channels.
  \item \textbf{Data Persistence Tier:} Managed PostgreSQL 15 relational database comprising 20 normalized tables protected by kernel-level Row-Level Security (RLS) policies and append-only audit logs.
\end{enumerate}

\begin{table}[htbp]
  \caption{AURORA Four-Tier Cloud Architecture Stack}
  \label{tab:architecture}
  \centering
  \footnotesize
  \begin{tabular}{p{2.1cm}p{2.4cm}p{3.1cm}}
    \toprule
    \textbf{System Tier} & \textbf{Core Stack} & \textbf{Security \& Governance Role} \\
    \midrule
    1. Client UI & React 19, Tailwind v4, \texttt{PDF.js} & Normalized coordinate pinning; SHA-256 backup verifier \\
    2. Edge / API & Next.js 16, Server Actions & RBAC middleware; RA 8792 step-up auth; conflict guard \\
    3. BaaS Services & Supabase Auth, S3, WSS & Signed private URLs; live WAL WebSocket sync \\
    4. Persistence & PostgreSQL 15, RLS, WAL & 20-table RLS isolation; atomic serial RPC; PITR backups \\
    \bottomrule
  \end{tabular}
\end{table}

\subsection{Core Algorithmic & Security Mechanisms}

\subsubsection{Normalized Coordinate PDF Annotation}
When an adviser or panelist highlights a manuscript passage on page $p$, the client transforms screen pixel offsets $(X_{\text{px}}, Y_{\text{px}}, W_{\text{px}}, H_{\text{px}})$ relative to the rendered page canvas dimensions $(C_W, C_H)$ into resolution-independent percentages:
\begin{equation}
  x\% = \frac{X_{\text{px}}}{C_W} \times 100, \quad
  y\% = \frac{Y_{\text{px}}}{C_H} \times 100
\end{equation}
This guarantees zero positional drift across zoom levels ($75\%$--$200\%$).

\subsubsection{Weighted Rubric Engine & RA 8792 Digital Locking}
For a rubric template containing $n$ criteria with normalized weights $\sum_{i=1}^{n} w_i = 100\%$, the evaluator's composite weighted score $S_{\text{weighted}}$ is computed as:
\begin{equation}
  S_{\text{weighted}} = \sum_{i=1}^{n} \frac{r_i \cdot w_i}{100}
\end{equation}
where $r_i \in [0, 100]$ is the raw criterion score. Upon clicking \textit{Sign \& Finalize}, the server re-authenticates the panelist's password, invokes the atomic PostgreSQL sequence \texttt{generate\_certificate\_serial()} to issue serial $\mathcal{S} = \texttt{AURORA-YYYY-NNNNNN}$, and computes a deterministic SHA-256 digest:
\begin{equation}
  H_{\text{cert}} = \text{SHA-256}\big(\text{evalId} \;\|\; \text{verdict} \;\|\; \mathcal{S} \;\|\; t_{\text{signed}}\big)
\end{equation}
Once committed with $\text{status} = \texttt{submitted}$, a PostgreSQL RLS policy blocks any subsequent \texttt{UPDATE} or \texttt{DELETE} operations on the record.

\subsection{Instrument Validation & Empirical Evaluation Setup}
To evaluate Objective 4, a structured survey instrument combining the 10-item \textbf{System Usability Scale (SUS)} \cite{brooke1996sus} and the \textbf{Unified Theory of Acceptance and Use of Technology (UTAUT)} \cite{venkatesh2003user} was subjected to expert content validation by three institutional validators using the Good and Scates (1972) criteria, achieving a validity rating of \textbf{4.89 / 5.00 (Very Highly Valid)}. Following expert refinement, two redundant items (original \texttt{PE2} and \texttt{UB2}) were excluded and subsequent items were renumbered across six UTAUT constructs: Performance Expectancy (\texttt{PE1--PE3}), Effort Expectancy (\texttt{EE1--EE4}), Social Influence (\texttt{SI1--SI4}), Facilitating Conditions (\texttt{FC1--FC4}), Behavioral Intention (\texttt{BI1--BI3}), and Use Behavior (\texttt{UB1--UB9}).

Internal consistency reliability across all 37 items (10 SUS + 27 UTAUT) was verified using Cronbach's Alpha across $N = 27$ respondents (15 BSIT senior students, 5 research advisers, 5 defense panelists, 1 capstone coordinator, and 1 college dean):
\begin{equation}
  \alpha = \frac{K}{K - 1}\left(1 - \frac{\sum_{i=1}^{K} \sigma_i^2}{\sigma_X^2}\right) = \mathbf{0.946}
\end{equation}
confirming excellent psychometric reliability ($\alpha > 0.70$).

%% ============================================================================
%% 4. RESULTS AND DISCUSSION
%% ============================================================================
\section{Results and Discussion}

\subsection{Workflow Efficiency & Turnaround Time Reduction}
Deployment of AURORA across simulated and live capstone defense workflows demonstrated substantial operational improvements over the legacy paper-based process (Table~\ref{tab:efficiency}). By replacing physical manuscript printing and manual routing with SHA-256 versioned uploads and instant Form DCS-CF-03 digital endorsement, average adviser review turnaround time decreased from \textbf{18.4 days to 2.6 days}---an \textbf{85.9\% reduction}. Furthermore, automated room-clash and adviser conflict-of-interest validation eliminated scheduling collisions entirely (0 conflicts across 42 test executions).

\begin{table}[htbp]
  \caption{Operational Comparison: Traditional Paper Process vs. AURORA Platform}
  \label{tab:efficiency}
  \centering
  \footnotesize
  \begin{tabular}{p{3.2cm}cc}
    \toprule
    \textbf{Operational Metric} & \textbf{Paper Workflow} & \textbf{AURORA Platform} \\
    \midrule
    Adviser Endorsement Turnaround & 18.4 days & \textbf{2.6 days ($-85.9\%$)} \\
    Paper Printed per Group (4 Stages) & 1,800+ pages & \textbf{0 pages (100\% Paperless)} \\
    Rubric Consolidation & 45--60 mins & \textbf{$< 1.0$ sec (Instant)} \\
    Room / Panelist Scheduling Clashes & 14.2\% rate & \textbf{0.0\% (Hard Guard)} \\
    Disaster Recovery Backup RTO / RPO & None (Manual) & \textbf{$\le 4\text{h}$ RTO / WAL PITR} \\
    \bottomrule
  \end{tabular}
\end{table}

\subsection{Quality Assurance, Security & Disaster Recovery Verification}
A comprehensive suite of 42 functional, security, and integration test cases (Annex H) was executed against the staging environment, achieving a \textbf{100.0\% pass rate (42/42 Passed)}. Penetration tests confirmed that:
\begin{itemize}
  \item Cross-tenant queries by unauthorized students return zero rows under PostgreSQL RLS (\texttt{TC-02}).
  \item Attempting to assign a project's Research Adviser to its evaluation panel is deterministically blocked by the Conflict-of-Interest Firewall (\texttt{TC-05}).
  \item Attempting to mutate a signed evaluation record fails at the database kernel level (\texttt{TC-09}).
  \item The integrated \textbf{Maintenance, Troubleshooting \& Disaster Recovery Console} (\texttt{/admin/system-health}) successfully executes 6 live subsystem probes, repairs orphaned version pointers via one-click self-healing maintenance, and exports cryptographically signed 20-table JSON cold archives verified via client-side SHA-256 checksums.
\end{itemize}

\subsection{Objective 4: Empirical SUS & UTAUT Evaluation Results}
Table~\ref{tab:utaut_results} summarizes the empirical evaluation gathered from $N = 27$ ParSU CECS stakeholders. On the 10-item System Usability Scale (SUS), AURORA achieved a composite score of \textbf{91.8 out of 100} (normalized mean $\mu = 4.84, \sigma = 0.28$), corresponding to an \textbf{``A'' Grade (Excellent / Best Imaginable Usability)}.

\begin{table}[htbp]
  \caption{Empirical Usability (SUS) and Technology Acceptance (UTAUT) Results ($N = 27$)}
  \label{tab:utaut_results}
  \centering
  \footnotesize
  \begin{tabular}{p{3.9cm}ccc}
    \toprule
    \textbf{Construct / Evaluation Dimension} & \textbf{Mean ($\mu$)} & \textbf{SD ($\sigma$)} & \textbf{Interpretation} \\
    \midrule
    Part I: System Usability Scale (SUS1--10) & 4.84 & 0.28 & Grade A (91.8/100) \\
    Part II.1: Performance Expectancy (PE1--3) & 4.88 & 0.22 & Highly Acceptable \\
    Part II.2: Effort Expectancy (EE1--4) & 4.86 & 0.25 & Highly Acceptable \\
    Part II.3: Social Influence (SI1--4) & 4.81 & 0.30 & Highly Acceptable \\
    Part II.4: Facilitating Conditions (FC1--4) & 4.85 & 0.26 & Highly Acceptable \\
    Part II.5: Behavioral Intention (BI1--3) & 4.90 & 0.20 & Highly Acceptable \\
    Part II.6: Use Behavior (UB1--9) & 4.87 & 0.24 & Highly Acceptable \\
    \midrule
    \textbf{Composite Grand Mean} & \textbf{4.86} & \textbf{0.25} & \textbf{Highly Acceptable} \\
    \bottomrule
  \end{tabular}
\end{table}

Across the six UTAUT dimensions, Behavioral Intention achieved the highest rating ($\mu = 4.90, \sigma = 0.20$), followed by Performance Expectancy ($\mu = 4.88, \sigma = 0.22$) and Use Behavior ($\mu = 4.87, \sigma = 0.24$). Specifically, 100\% of student and faculty respondents strongly agreed that they prefer using AURORA over traditional paper-based defense processes (\texttt{UB9}) and intend to adopt the platform for future academic defenses (\texttt{BI1--BI3}).

%% ============================================================================
%% 5. CONCLUSION AND RECOMMENDATIONS
%% ============================================================================
\section{Conclusion and Recommendations}

This study successfully developed and validated \textbf{AURORA}, an end-to-end paperless academic defense management and quality assurance platform tailored to the institutional requirements of Partido State University. By unifying viewport-agnostic PDF coordinate annotations, conflict-free defense scheduling, dynamic weighted rubric grading, Republic Act No. 8792 step-up digital signatures, and a 3-2-1 disaster recovery console, AURORA eliminates the delays, printing expenses, and integrity risks of manual paper workflows. Empirical evaluation via SUS (91.8/100) and UTAUT ($\mu = 4.86/5.00$) confirms high usability and strong institutional readiness for university-wide adoption.

Future work includes integrating automated OCR similarity pre-screening directly within the PDF upload pipeline and expanding mobile push notifications via Progressive Web App (PWA) service workers.

\begin{acks}
The proponents express their sincere gratitude to their Research Adviser, \textbf{Prof. Froiland B. Pe\~{n}alosa}, Capstone Coordinator and Department Chairperson \textbf{Dr. Kennedy C. Cuya}, and College Dean \textbf{Engr. Emelina R. Padayao} of the College of Engineering and Computational Sciences, Partido State University, for their technical guidance and institutional support.
\end{acks}

\bibliographystyle{ACM-Reference-Format}
\bibliography{references}

\end{document}
